\lysh{13312}{4}

\begin{alist}
\item Η εξίσωση $(1)$ είναι $2$ου βαθμού με διακρίνουσα $\Delta = 36 - 4\lambda$ και για να έχει πραγματικές ρίζες πρέπει και αρκεί $\Delta \ge 0$ ισοδύναμα $36 - 4\lambda \ge 0$ οπότε $-4\lambda \ge -36$ και τελικά $\lambda \le 9$.

\item 
\begin{rlist}
\item Αφού οι πραγματικοί αριθμοί $\alpha, \beta$ έχουν σταθερό άθροισμα $6$ και γινόμενο $\lambda$ θα
είναι ρίζες της $(1)$ και αυτό όπως δείξαμε στο $\alpha)$ ερώτημα συμβαίνει αν και μόνο αν
\[\lambda \le 9 \Leftrightarrow \alpha \cdot \beta \le 9.\]
\item Αφού οι πραγματικοί αριθμοί $\alpha, \beta$ είναι ρίζες της $(1)$, θα είναι $\alpha = \beta$ αν και μόνο αν
η $(1)$ έχει δύο ρίζες ίσες, δηλαδή ισοδύναμα αν $\Delta = 0$ δηλαδή $36 - 4\lambda = 0$ οπότε
$-4\lambda = -36$ και τελικά $\lambda = 9$ πράγμα που σημαίνει ότι $\alpha \cdot \beta = 9$.
\end{rlist}

\item Έστω $\alpha, \beta$ οι διαστάσεις τυχαίου ορθογωνίου παραλληλογράμμου με περίμετρο $12$.
Τότε $2\alpha + 2\beta = 12$ δηλαδή $\alpha + \beta = 6$. Το εμβαδόν τους είναι ίσο με $\alpha \cdot \beta$.
Δείξαμε στο $\beta)$ ερώτημα ότι αν για δύο πραγματικούς αριθμούς $\alpha, \beta$ ισχύει ότι $\alpha + \beta = 6$,
τότε $\alpha \cdot \beta \le 9$ και μάλιστα $\alpha \cdot \beta = 9$ αν και μόνο αν $\alpha = \beta$.

Συνεπώς η μεγαλύτερη τιμή του εμβαδού είναι $9$ και αυτό συμβαίνει αν και μόνο αν οι
διαστάσεις του ορθογωνίου είναι ίσες, δηλαδή αν και μόνο αν είναι τετράγωνο.
\end{alist}
